\documentclass[11pt,a4paper]{article}
\usepackage{cmap}
\usepackage[T1]{fontenc}
\usepackage[utf8]{inputenc}
\usepackage[spanish,es-nodecimaldot,es-noshorthands]{babel}
\usepackage[margin=24mm]{geometry}
\usepackage{booktabs}
\usepackage{pgfplots}
\usepackage[hidelinks]{hyperref}
\input{glyphtounicode}
\pdfgentounicode=1
\pgfplotsset{compat=1.18,
  barras/.style={xbar,bar width=8pt,width=0.87\linewidth,height=4.5cm,
    xmin=0,axis lines=left,xmajorgrids,ytick=data,y dir=reverse,
    enlarge y limits=0.25,tick label style={font=\small},
    fill=blue!55,draw=blue!65!black},
  /pgf/number format/use comma,
  /pgf/number format/1000 sep={\,}}
\newcommand{\codigo}[1]{\texttt{#1}}
\newcommand{\fuente}[1]{\par\smallskip\noindent\textit{Fuente:} #1\par}

\hypersetup{pdftitle={Etapa 3: Búsquedas y agregaciones en Elasticsearch},
  pdfauthor={Camilo Fuentes, Pablo Adasme, María Ignacia Parra Muñoz}}

\begin{document}
\begin{center}
{\LARGE\bfseries Etapa 3: Búsquedas y agregaciones en Elasticsearch\par}
\smallskip
Camilo Fuentes \quad Pablo Adasme \quad María Ignacia Parra Muñoz
\end{center}

\section*{Resumen}
El orden de búsqueda cambió la composición de los diez primeros pedidos sin
cambiar el total de coincidencias. En vuelos, el límite bajo de conteo devolvió
una cota inferior y no un total exacto. En registros web, las cuatro estrategias
coincidieron en los promedios redondeados, pero no en la latencia. Estos resultados
corresponden a tres índices de ejemplo distintos en una instancia local; no
constituyen una prueba sobre un conjunto común ni sobre 1~GB de datos o diez
millones de documentos.

\section{Entorno y reproducción}
Se utilizó Docker Compose con Elasticsearch y Kibana 9.5.6 y JupyterLab.
Elasticsearch se ejecutó en un solo nodo, con un límite de 1~GB de RAM para
su contenedor. Este límite no es el tamaño del conjunto de datos.

\begin{enumerate}
\setlength{\itemsep}{0pt}
\item Con Docker en ejecución, acceso a Internet para la primera descarga y
los puertos 9200, 5601 y 8888 libres, ejecutar desde la raíz del repositorio:
\begin{verbatim}
docker compose up -d --build
docker compose ps
\end{verbatim}
\item Abrir \href{http://localhost:5601}{Kibana}, esperar a que esté listo
y buscar \textit{Integrations} y luego \textit{Sample Data}. Instalar los tres
conjuntos oficiales si no están cargados. Consultar la
\href{https://www.elastic.co/docs/manage-data/ingest/sample-data}{guía oficial}
si cambia la interfaz.
\item En \textit{Dev Tools} $\to$ \textit{Console}, ejecutar por separado:
\begin{verbatim}
GET kibana_sample_data_ecommerce/_count
GET kibana_sample_data_flights/_count
GET kibana_sample_data_logs/_count
\end{verbatim}
Comprobar \codigo{count} y que \codigo{\_shards.failed} sea 0.
\item Abrir el
\href{http://localhost:8888/lab/tree/work/00_indice.ipynb}{cuaderno índice}
en JupyterLab. Ejecutar primero las celdas de conexión y configuración y
luego las mediciones y su comparación. El kernel en Compose usa
\url{http://elasticsearch:9200}; desde el equipo anfitrión se usa
\url{http://localhost:9200}. Los servicios solo están publicados en la
interfaz local. Elasticsearch y Jupyter no tienen autenticación en esta
demostración: no deben exponerse a la red.
\end{enumerate}

\begin{center}
\begin{tabular}{llr}
\toprule
Conjunto & Índice de Elasticsearch & Documentos\\
\midrule
Pedidos & \codigo{kibana\_sample\_data\_ecommerce} & 4\,675\\
Vuelos & \codigo{kibana\_sample\_data\_flights} & 13\,014\\
Registros web & \codigo{kibana\_sample\_data\_logs} & 14\,074\\
\bottomrule
\end{tabular}
\end{center}
Los conteos sin filtro devolvieron un fragmento correcto y ninguno fallido
por índice. Son conteos locales, no tamaños en disco ni valores garantizados
para otras instalaciones. Se eligieron muestras oficiales y pequeñas para
repetir las pruebas en el entorno local y verificar los resultados sin
preparar una carga masiva. Cada muestra aporta campos adecuados para su
pregunta; eso no demuestra representatividad a mayor escala.

\paragraph{Demostración en vivo.}
Conviene iniciar los servicios y comprobar los conteos antes de presentar.
En los experimentos 1 y 2, copiar en Kibana la petición completa de cada
medición; contrastar, respectivamente, \codigo{hits.hits} y
\codigo{hits.total} con la tabla de Jupyter. En el 3, contrastar
\codigo{aggregations} y ejecutar la sección 5 del cuaderno para medir las
estrategias. Mantener los mismos parámetros en ambas interfaces. Si la
instancia no está disponible, presentar las tablas del informe como
resultados locales anteriores, no como una ejecución en vivo.

\section{Experimento 1: ordenar cambia los primeros diez pedidos}
\textbf{Autor:} Camilo Fuentes. \textbf{Datos:} pedidos de comercio electrónico.

\paragraph{Pregunta e hipótesis.}
¿Cómo cambia el importe medio con impuestos de los primeros diez pedidos en
euros que coinciden con \codigo{clothing} al variar el orden? Se esperaba que
la media por relevancia quedara entre las obtenidas por importe ascendente
y descendente.

\paragraph{Diseño.}
La variable independiente fue el criterio principal de orden:
\codigo{taxful\_total\_price} ascendente o descendente, o \codigo{\_score}
descendente. Se mantuvieron fijos el índice, el \codigo{match} de
\codigo{clothing} en \codigo{category}, el filtro exacto \codigo{currency: EUR},
\codigo{size: 10} y el desempate por \codigo{order\_id} ascendente.
La variable dependiente fue la media aritmética de
\codigo{taxful\_total\_price} de los diez pedidos devueltos, con los importes
originales de \codigo{\_source}, no los valores indexados usados para ordenar.
El cuaderno comprueba que el total sea exacto y que haya diez importes válidos
en EUR; no publica la comparación con resultados parciales.

\begin{center}
\begin{tabular}{lrrr}
\toprule
Criterio principal & Coincidencias & Pedidos observados & Media (EUR)\\
\midrule
Importe ascendente  & 3\,927 & 10 & 15,981\\
Importe descendente & 3\,927 & 10 & 526,456\\
Relevancia          & 3\,927 & 10 & 55,980\\
\bottomrule
\end{tabular}
\end{center}

\begin{center}
\begin{tikzpicture}
\begin{axis}[barras,xmax=600,ymin=-0.75,ymax=2.75,
  yticklabels={Ascendente,Descendente,Relevancia},
  xlabel={Importe medio con impuestos (EUR)}]
\addplot coordinates {(15.981,0) (526.456,1) (55.980,2)};
\end{axis}
\end{tikzpicture}
\end{center}

\paragraph{Conclusión.}
La relevancia quedó entre ambos extremos y la diferencia entre descendente
y ascendente fue de 510,475~EUR. El orden cambió qué pedidos entraron entre
los primeros diez, no sus importes ni las 3\,927 coincidencias. Las medias
describen los diez resultados de cada orden, no todos los pedidos coincidentes
ni un efecto causal generalizable.
\fuente{\nolinkurl{jupyter/notebooks/01_busqueda_ordenamiento.ipynb}.}

\clearpage
\section{Experimento 2: un total puede ser una cota inferior}
\textbf{Autor:} Pablo Adasme. \textbf{Datos:} vuelos.

\paragraph{Pregunta.}
¿Cuándo es exacto el total de una búsqueda de vuelos hacia Australia?

\paragraph{Diseño.}
Solo cambió \codigo{track\_total\_hits}: 10, 1000 o \codigo{true}. Se usó el
mismo índice y se fijó \codigo{DestCountry: AU}, el rango de
\codigo{AvgTicketPrice} en $[500,950)$ (incluye 500 y excluye 950) y
\codigo{size: 0} para omitir los vuelos individuales. Se leyeron juntos
\codigo{hits.total.value} y \codigo{hits.total.relation}.
\codigo{eq} indica un total exacto; \codigo{gte}, un límite inferior.

\begin{center}
\begin{tabular}{lrrl}
\toprule
\codigo{track\_total\_hits} & Valor & Relación & Lectura\\
\midrule
\codigo{10}   & 10  & \codigo{gte} & Al menos 10\\
\codigo{1000} & 242 & \codigo{eq}  & Exactamente 242\\
\codigo{true} & 242 & \codigo{eq}  & Exactamente 242\\
\bottomrule
\end{tabular}
\end{center}

\begin{center}
\begin{tikzpicture}
\begin{axis}[width=0.78\linewidth,height=4.2cm,axis lines=left,xmajorgrids,
  xmin=0,xmax=280,ymin=-0.5,ymax=2.5,y dir=reverse,ytick={0,1,2},
  yticklabels={10,1000,true},ylabel={\codigo{track\_total\_hits}},
  xlabel={Coincidencias (valor devuelto)}]
\addplot[only marks,mark=*,blue!65!black] coordinates {(242,1) (242,2)};
\draw[->,thick,orange!80!black] (axis cs:10,0) -- (axis cs:36,0);
\addplot[only marks,mark=o,mark options={fill=white},orange!80!black]
  coordinates {(10,0)};
\end{axis}
\end{tikzpicture}
\par\small Círculo relleno: total exacto. Círculo vacío con flecha:
cota inferior; la longitud de la flecha no representa un conteo.
\end{center}

\paragraph{Conclusión.}
Hubo 242 coincidencias para el filtro fijo. El umbral 10 no acreditó el
total exacto; el de 1000 coincidió con la referencia \codigo{true}.
Esto no demuestra que el umbral 1000 sea suficiente para otros filtros
ni mide la latencia de las opciones.

\paragraph{Comprobación adicional del cuaderno.}
El anexo A realiza otras tres consultas con \codigo{/\_count}: mantiene
\codigo{AU} y cambia el rango de precio. Dio 70 vuelos en $[500,650)$,
76 en $[650,800)$ y 96 en $[800,950)$. Los intervalos no se solapan y
suman 242; el último tuvo 20 vuelos más que el siguiente por conteo.
Este anexo no forma parte de la comparación de umbrales.
\fuente{\nolinkurl{jupyter/notebooks/02_experimento.ipynb}.}

\clearpage
\section{Experimento 3: estrategias para agrupar y promediar}
\textbf{Autora:} María Ignacia Parra Muñoz. \textbf{Datos:} registros web.

\paragraph{Pregunta e hipótesis.}
Al agrupar registros web por código de respuesta HTTP (200, 404, 503)
y calcular el promedio de bytes transferidos, ¿cuánto cambia el tiempo de
respuesta según la forma interna en que Elasticsearch ejecuta el cálculo?
Se esperaba que la caché obtuviera la menor latencia al repetir la misma
consulta y que, sin caché, leer columnas preparadas (\codigo{doc\_values})
fuera más rápido que calcular el código al vuelo leyendo el JSON original.

\paragraph{Diseño.}
La variable independiente fue la estrategia de ejecución:
\begin{enumerate}
\item \textbf{Caché activa} (\codigo{request\_cache=true}): reutiliza el
resultado de la petición repetida a nivel de fragmento.
\item \textbf{Estándar sin caché} (\codigo{global\_ordinals}): lee
\codigo{doc\_values} y agrupa con identificadores enteros de los códigos.
\item \textbf{Tabla hash sin caché} (\codigo{execution\_hint: "map"}): lee
\codigo{doc\_values} y agrupa directamente los valores en una tabla hash.
\item \textbf{Cálculo al vuelo sin caché} (\codigo{runtime\_mappings}):
crea un campo temporal para el código HTTP leyendo \codigo{\_source}
y ejecutando un script Painless por registro.
\end{enumerate}
Se mantuvieron fijos el índice, la agrupación \codigo{terms} por código HTTP
(hasta cinco grupos), el promedio \codigo{avg} de \codigo{bytes},
\codigo{size: 0} y el mismo contenedor de un nodo con 1~GB de RAM.

\paragraph{Medición.}
Se hicieron 20 repeticiones tras tres ejecuciones de calentamiento por
configuración. Se registraron la latencia total mediana desde Python y el
\codigo{took} medio de Elasticsearch. En otras 20 peticiones por estrategia
sin caché se registró la mediana de la suma de tiempos de agregación
reportados por la API Profile. Es una medición instrumentada, no un tiempo
puro del motor ni un componente que deba restarse de la latencia total.
No se perfiló la configuración con caché.

\paragraph{Resultados.}
Las cuatro configuraciones devolvieron los mismos códigos y promedios al
redondear a dos decimales: 5\,897,85 bytes (200), 5\,049,24 bytes (404)
y 0,00 bytes (503).
\begin{center}
\begin{tabular}{lrrr}
\toprule
Configuración & Mediana total (ms) & \codigo{took} medio (ms) & Profile ($\mu$s)\\
\midrule
Caché activa  & 3,28  & 1,05  & No aplica\\
Estándar      & 3,72  & 1,05  & 1\,145,6\\
Tabla hash    & 5,72  & 2,50  & 2\,074,2\\
Al vuelo      & 94,02 & 91,75 & 93\,962,4\\
\bottomrule
\end{tabular}
\end{center}

\clearpage
\subsection*{Comparación de latencia y perfil}
\begin{center}
\begin{tikzpicture}
\begin{axis}[barras,xmax=100,ymin=-0.75,ymax=3.75,
  yticklabels={Caché,Estándar,Hash,Al vuelo},
  xlabel={Latencia total mediana (ms)}]
\addplot coordinates {(3.28,0) (3.72,1) (5.72,2) (94.02,3)};
\end{axis}
\end{tikzpicture}

\begin{tikzpicture}
\begin{axis}[barras,xmax=100000,ymin=-0.75,ymax=2.75,scaled x ticks=false,
  xticklabel style={/pgf/number format/fixed},
  xtick={0,25000,50000,75000,100000},
  yticklabels={Estándar,Hash,Al vuelo},
  xlabel={Mediana de agregación instrumentada ($\mu$s)}]
\addplot coordinates {(1145.6,0) (2074.2,1) (93962.4,2)};
\end{axis}
\end{tikzpicture}
\end{center}

\paragraph{Conclusión.}
En esta prueba, el cálculo al vuelo tardó 94,02~ms en total y registró
unas 82 veces el tiempo de agregación perfilado del estándar. En esta
implementación, leer \codigo{\_source} y ejecutar Painless por registro
añade trabajo a cambio de poder definir grupos sin reindexar. La tabla
hash registró unas 1,8 veces el tiempo perfilado del estándar: ambos
leyeron \codigo{doc\_values}, pero el estándar agrupó mediante identificadores
enteros en lugar de usar los valores de texto como claves.
La caché obtuvo la menor mediana total (3,28~ms), frente a 3,72~ms del
estándar y casi 29 veces menos que los 94,02~ms del cálculo al vuelo.
El \codigo{took} medio fue el mismo con caché y con el método estándar;
la medición desde Python incluye sobrecarga fuera del motor, pero estas
mediciones no aíslan el costo del viaje HTTP.
\fuente{\nolinkurl{informe.md} y sección 5 de
\nolinkurl{jupyter/notebooks/03_experimento.ipynb}. Gráficos reproducidos
con los valores reportados, no con una nueva medición.}

\section{Escalabilidad: plan, no resultado}
Como hipótesis para registros web, el informe original proyecta
\emph{más de un minuto} para el cálculo al vuelo con diez millones de
registros en un nodo. Es una extrapolación sin medir: multiplicar
aproximadamente 94~ms por 700 supone un crecimiento lineal bajo condiciones
constantes. No se verificó esa latencia ni el tamaño en disco a ese volumen.

Para probarlo, se propone aumentar el volumen por etapas y comparar las
mismas agregaciones con y sin un filtro temporal acotado, con
\codigo{doc\_values} y con campo al vuelo, separando caché fría de caliente.
Después se podría distribuir el índice en varios fragmentos primarios
en distintos nodos para comparar el trabajo local y la reducción final.
Registrar volumen real, tamaño en disco, configuración y latencias
permitiría verificar si el paralelismo y la caché de consultas repetitivas
aportan las mejoras previstas.

\end{document}
